\section{学位论文研究方案}

\subsection{研究目标}

本研究围绕弱连接交往对大学生价值观的影响及其应对策略，拟实现以下三个目标。

第一，明确大学生弱连接交往的范围与特征，构建考察其价值影响的分析框架。结合马克思主义经典著作、思想政治教育学研究和相关政策文献，明确弱连接交往的识别依据，
并从个人生活、人际交往、学习劳动与择业、社会公共四个方面梳理可观察的价值议题。

第二，考察弱连接交往对大学生价值观的具体影响。分析多平台公开网络文本，梳理学生在校园“搭子”、跨校学习互助和社群合作等情境中的交往叙述与价值表达，
考察材料中呈现的积极作用及可能带来的困惑。

第三，从思想政治教育的学科视角出发，提出贴合高校日常育人工作的引导策略。结合学生在具体交往情境中的需要与困惑，
从思政课教学、校园协作项目和社会实践等环节提出可操作的做法，为高校价值观教育提供参考。

\subsection{研究内容}

首先进行思想溯源与理论分析。以马克思关于人的本质与社会关系的论述为起点，梳理马克思主义关于社会交往、劳动与人的全面发展的相关观点，
并结合思想政治教育学对价值观形成与引导的阐释，说明从日常交往考察大学生价值观的理论依据\footcite{MarxEngels2009,sizheng2018}。

其次进行概念界定与框架构建。结合弱关系理论，从交往频率、情感投入、亲密程度与互惠情况辨识关系强度，将交往对象可辨认、存在实际回应并具有一定持续性作为识别具体交往的依据。
本文考察大学生对何者值得追求、不同目标如何取舍的基本看法。对行为善恶、责任和规范的判断，是其中涉及道德的部分。

《新时代公民道德建设实施纲要》将社会公德、职业道德、家庭美德和个人品德建设列为着力点，为识别不同生活场景中的规范议题提供参照\footcite{gangyao2019}。
《思想道德与法治》（2023年版）第一章第一节讨论价值观如何影响人生目的和人生道路的选择，
第五章第三节分述公共生活、职业生活、家庭生活中的道德要求和个人品德修养\footcite{sixiangdaode2023}。
《伦理学》（第二版）第八章第一节讨论一般价值与道德价值的关系，第六章第四节阐明社会主义核心价值观对社会主义道德规范的统领作用\footcite{lunlixue2020}。
习近平总书记在2014年5月4日北京大学师生座谈会上指出，“核心价值观，其实就是一种德”，并提出青年要在勤学、修德、明辨、笃实上下功夫，
为理解青年如何在日常实践中践行社会主义核心价值观提供了依据\footcite{Xi2014}。

结合上述论述，本文按照大学生的实践领域和价值观所指向的对象，构建四维分析框架。个人生活价值观关注消费选择、闲暇安排和精神追求。
人际交往价值观关注大学生与同伴、家人等具体他人交往时如何理解边界、互惠与承诺。学习劳动与择业价值观关注专业学习的意义、劳动价值及职业选择标准。
社会公共价值观关注面向公共生活的规则意识、参与意愿和社会责任。

随后考察现实表现与价值影响。以多平台公开主题帖和关联回复为材料，归纳大学生弱连接交往的常见情境与讨论主题，比较个人生活、人际交往、学习劳动与择业、社会公共四个方面的价值表达。
结合帖文中的经历叙述与互动上下文，考察信息互通与合作互助呈现的积极作用，以及片面比较、短期利益取向和公共责任意识弱化等可能出现的问题，辨析其具体表现与适用情境。

最后提出思想政治教育应对策略。针对网络文本分析梳理出的具体问题，探索如何将学生关心的交往议题引入思政课堂，如何借助校园协作项目明确分工、约定与责任，
以及如何通过志愿服务和社会实践发挥跨群体联系的育人作用。根据场景和学生需求说明各项做法的适用范围。

\subsection{拟采用的研究方法与技术路线}

\subsubsection{研究方法}

本研究坚持辩证唯物主义和历史唯物主义的根本方法。文献研究用于确定理论依据和分析框架，网络文本挖掘用于梳理交往情境及讨论主题，质性内容分析用于结合语境解读价值表达。
三种方法依次回答研究什么、材料呈现了什么以及如何理解这些表现。

\paragraph{文献研究法}
研读马克思主义经典著作、党的创新理论文献、高校思想政治教育政策和相关教材，梳理社会交往与青年价值观形成的理论依据。
结合国内外弱关系、大学生交往和价值观研究，明确弱连接交往的识别标准与四维分析框架，据此确定语料检索范围和初步编码维度。

\paragraph{网络文本挖掘法}
拟从高校公开论坛、小红书、豆瓣高校及青年小组、知乎高等教育及应届生板块等来源，按平台收集近三年与大学生弱连接交往有关的公开主题帖及回复。
结合板块属性和发帖者公开自述等线索判断学生身份，对身份或交往情境不清楚的材料予以排除。
语料围绕校园“搭子”、跨校备考求职互助和兴趣社群等场景选取，保留交往对象可辨认、有实际回应且联系具有一定延续性的叙述和讨论，剔除单向浏览、一次性争论等材料。
按平台公开访问规则，结合人工检索与Python辅助整理进行去重、去标识化和人工复核，保留主题帖与关联回复的上下文。

对整理后的语料进行中文分词、关键词提取、共现分析和短文本主题聚类，查看各平台中常见的交往情境、讨论主题及其与价值议题的关联。
将用于选取材料的场景词与待分析的价值表达区分，避免把检索词的高频出现直接解释为研究发现。依据主题分布和材料完整程度，为下一步质性分析确定抽样范围。

\paragraph{质性内容分析法}
结合挖掘得到的主题分布，从不同平台选取典型主题帖及关联回复，并纳入表达存在差异或与常见模式不一致的材料。
依据四维框架编写编码说明，借助NVivo进行人工编码，通过试编码、分歧复核和跨平台对照校正判断。
分析时结合发言上下文识别调侃、反讽和意见分歧，辨析学生如何评价交往经历及其与价值观的关系，并将质性结果与文本挖掘归纳出的主题相互印证。
研究结论以公开文本中可观察的表达及其语境为依据，解释范围限于样本所涉交往场景。

\subsubsection{技术路线}

研究先确定理论依据和分析框架，再通过网络文本挖掘归纳讨论主题，结合质性内容分析解释价值表达，最后提出育人策略。具体流程见图\ref{fig:research-roadmap}。

\begin{figure}[htbp]
  \centering
  \small
  \renewcommand{\thefigure}{\arabic{figure}}
  \setlength{\fboxsep}{5pt}
  \begin{tabular}{@{}c@{}}
    \fbox{\parbox{0.8\linewidth}{\centering\textbf{思想溯源与理论分析}\\梳理社会交往与青年价值观形成的理论依据}} \\
    \(\downarrow\) \\
    \fbox{\parbox{0.8\linewidth}{\centering\textbf{概念界定与框架构建}\\明确弱连接交往的识别依据与价值观分析维度}} \\
    \(\downarrow\) \\
    \fbox{\parbox{0.8\linewidth}{\centering\textbf{网络文本收集与主题挖掘}\\筛选公开帖文和回复，归纳交往场景与讨论主题}} \\
    \(\downarrow\) \\
    \fbox{\parbox{0.8\linewidth}{\centering\textbf{质性编码与价值影响分析}\\结合语境解读价值表达，辨析积极作用与现实困惑}} \\
    \(\downarrow\) \\
    \fbox{\parbox{0.8\linewidth}{\centering\textbf{提出思想政治教育应对策略}\\面向思政课、校园协作和社会实践提出具体做法}}
  \end{tabular}
  \caption{研究技术路线图}
  \label{fig:research-roadmap}
\end{figure}

\subsection{研究的创新点与难点}

\subsubsection{研究可能存在的创新点}

\paragraph{研究方法}
本研究在同一批多平台网络文本中衔接主题挖掘与质性内容分析。先通过主题挖掘掌握各类交往议题在语料中的分布，再沿主题回到原帖及关联回复，选取典型和意见分歧较大的文本进行编码。
比较同一主题下学生表达的立场及其理由，可以说明样本中哪些议题较常出现，以及这些议题为何会得到不同理解，并使分析结论能够追溯到具体语料。

\paragraph{研究内容}
已有研究分别讨论了弱连接交往的信息互通、求职支持、交往边界和社会比较。本研究在个人生活、人际交往、学习劳动与择业、社会公共四个维度下比较校园“搭子”、跨校互助和兴趣社群等不同情境，
考察同一种交往实践可能涉及哪些价值议题。例如，跨校求职互助既可能扩展学生对职业选择的认识，也可能使其更频繁地把录用结果用作评价个人发展的参照。
通过比较具体情境，呈现弱连接交往的积极作用与学生遇到的价值困惑各有哪些表现、在何种情境中同时出现。

\paragraph{研究视角}
本研究从马克思主义关于人的本质与社会关系的论述出发，把大学生的价值观放在具体的相处关系中理解，关注交往双方如何理解彼此的承诺与责任。
例如，任务结束后减少往来，学生可能将其理解为尊重边界，也可能在对方没有履行既有约定时将其视为失信。
分析学生作出判断时援引的理由，可以说明个人边界、互惠与责任在实际交往中如何被理解，为高校价值观教育找到具体的讨论议题。

\subsubsection{研究难点}

一是将四维框架转化为可一致执行的编码规则。一则兴趣社群帖文可能同时谈到个人兴趣、合作承诺和群体规则，编码时需要确定每段话表达的主要判断及其涉及的维度。
研究拟以相对完整的价值观表述为编码单位，保留原帖及回复供复核，通过试编码完善各维度的纳入标准。允许一条文本涉及多个维度，汇总时分别报告文本条数与编码次数，避免重复计数影响比较。

二是准确运用经典著作和相关历史文献。有关社会关系、劳动与人的全面发展的论述形成于特定历史语境，转用于解释当代大学生的交往实践时，需要说明其适用范围。
研究拟核对原著中的概念、论述对象和上下文，交代所援引观点能够阐明的问题。分析中分别说明经典论述提供的理论依据、当代政策文献提出的育人要求，以及网络文本能够支持的经验判断。

三是保证多平台语料的可比性，并控制结论的解释范围。各平台的发帖形式、回复结构及可检索内容存在差异，互动量也可能受平台传播方式影响。
研究拟统一记录来源、发布时间和交往情境，采用一致的纳入与排除标准，按平台和主题复核语料覆盖情况，适时补充有代表性的新情境。
分析时交代身份线索和材料分布的局限，将结论限定在样本能够支持的价值表达及其情境。
